\documentclass[12pt]{article}
\usepackage[utf8]{inputenc}
\usepackage[T1]{fontenc}
\usepackage{hyperref}
\title{Anmatyerr Language Resources for K-12 Curriculum Development}
\author{Compiled by Claude Sonnet 5.5 (claude-sonnet-5-5) for Andr\'as Kornai}
\date{}
\begin{document}
\maketitle

\section{Introduction}
Anmatyerr (ISO 639-3 amx; Glottocode anma1239; also spelled Anmatjirra, Anmatjera, Unmatjera) is a Pama-Nyungan language of Central Australia, spoken mainly in communities about 200 km north of Alice Springs (for example Ti Tree, Laramba and Mount Allan). AIATSIS lists a Central and an Eastern dialect. This report is based on a small number of web searches; it found no school curriculum programme for Anmatyerr, and none is claimed here.

\section{Primary Educational Resources}
\begin{itemize}
\item \emph{Central Anmatyerr Picture Dictionary} (Jenny Green, IAD Press, 2003), described by the publisher as an educational resource for school children, teachers and speakers improving literacy, with about 400 illustrations and definitions in Anmatyerr and English. I did not see a direct product URL for this title in the search results.
\end{itemize}

\section{Linguistic and Reference Materials}
\begin{itemize}
\item \emph{Central \& Eastern Anmatyerr to English Dictionary} (Jenny Green, IAD Press, 2010), a comprehensive dictionary: \url{https://www.iadpress.com.au/products/central-eastern-anmatyere-to-eng-dictionary}
\item AustLit record for related Anmatyerr publishing: \url{https://www.austlit.edu.au/austlit/page/C790542}
\item AIATSIS AUSTLANG entry for Anmatyerr (C8.1), listing dialects and further publications: \url{https://collection.aiatsis.gov.au/austlang/language/c8.1}
\end{itemize}

\section{Religious and Community Resources}
No religious or community-organisation resources were found in the searches.

\section{Digital and Audio Resources}
\begin{itemize}
\item Endangered Languages Project sample page for Anmatyerr: \url{https://endangeredlanguages.com/lang/3840/samples/6882}
\end{itemize}

\section{Curriculum Development Context}
Published materials include a school-oriented picture dictionary and a large bilingual dictionary, both by Jenny Green with Anmatyerr speakers. AIATSIS also lists local-authored titles, such as \emph{Anmatyerr Angkety Angkwey} (Margaret Carew, 2010) and \emph{Alherramp-areny angkety arelh-kenh} (Green and Peltharr, 2000), which could serve as reading texts. Current availability, and any school or bilingual programme, should be checked directly with the communities, IAD Press and the Northern Territory education authorities.
\end{document}